\section{Emergent relativity and effective Lorentz symmetry}
\label{sec:emergent-relativity}

The substrate carries Lorentzian sign structure only at the level of the brane \emph{action} (\Cref{subsec:ontology-kinematics}); the ambient $\R^4$ in which the brane is embedded is symmetric Euclidean, and on the spacelike-slice working description $t$ enters as the slicing parameter rather than as a metric direction.
A Lorentz \emph{metric} on excitations is therefore not postulated --- it is an emergent property of the linear wave regime as seen by an inside observer.
This section makes precise the sense in which Lorentz kinematics emerge as a kinematic symmetry of the dominant dispersion branch that controls information transport.

\subsection{Linear-wave cone and an effective light speed}

Linearizing \Cref{eq:eom} about a homogeneous background yields a multi-branch wave system.
For isotropic elastic media, one expects (at long wavelengths) transverse and longitudinal branches.
The characteristic speeds $c_L^2 = k_s a^2/m$ and $c_T^2 = (1-\alpha)\,k_s a^2/m$ are derived
in \Cref{sec:continuum-model} from two independent starting points --- the closed-form dispersion
of the discrete action (\Cref{eq:cone-speeds}) and the acoustic-tensor eigenvalues of the $W_2$
continuum limit (\Cref{eq:wave-speeds-derived}) --- and their agreement there is a self-consistency
check across levels of description.
For each polarization branch $b$ one has the long-wavelength dispersion
\begin{equation}
  \omega_b^2(\veck)\;\approx\; c_b^2\,|\veck|^2 \quad \text{for } |\veck|\to 0,
  \label{eq:linear-dispersion}
\end{equation}
which defines an effective wave cone for that branch.

In regimes where a single branch dominates transport and dispersion is negligible,
the wave equation is approximately invariant under the standard Lorentz group at speed $c$,
with boost parameter
\begin{equation}
  \gamma = \bigl(1-v^2/c^2\bigr)^{-1/2},
  \label{eq:lorentz}
\end{equation}
as an effective symmetry of the linear propagation sector, not as a microscopic postulate.

\paragraph{Dual-observer reading: the only Lorentz observer is the internal one.}
On the 6-neighbor axial-only lattice (\Cref{subsec:lattice-substrate}) the long-wavelength wave speed
is direction-dependent at the lab-frame level --- for example $c_L([100])/c_L([111])\approx 1.075$
at the default operating point $\alpha=0.2$.
The relativistic effective description \eqref{eq:lorentz} is therefore not a statement about the lab frame; it is a
conjecture about \emph{internal} observers, whose rods, clocks, and signal cones are substrate-emergent and renormalize
with the same direction-dependent speed (see \Cref{para:dual-observer} for the full mechanism).
Without observer universality the emergent Lorentz statement collapses to ``a particular branch
propagates with one speed along each direction'' and is not a symmetry of physics for any actual observer.

\subsection{Analogue-metric viewpoint}

An equivalent viewpoint, common in analogue-gravity systems, is that linear excitations experience an
effective metric structure with interval
\begin{equation}
  ds^2 = -c^2 dt^2 + g^{(\text{eff})}_{ij}(x,t)\,dx^i\,dx^j,
  \label{eq:effective-metric-interval}
\end{equation}
where $g^{(\text{eff})}_{ij}$ depends on the background configuration about which one linearizes.
In the present model there are two natural geometric objects: the induced brane metric $g_{ij}(\vecX)$
and the branch-dependent effective coefficients of the linearized operator.
In particular, slowly varying transverse profiles $u=X^4$ modify intrinsic distances already at the level of the induced metric (\Cref{eq:metric-near-identity,subsec:gravity-channel}), and can therefore be read as a background that focuses or defocuses wave transport in close analogy to a weak gravitational field.
Both encode how background deformations can bend or modulate wave propagation, in close analogy to
analogue-gravity constructions \citep{Barcelo2011AnalogueGravity,Volovik2003HeliumDroplet}.

\subsection{Regime of validity and expected deviations}

Effective Lorentz symmetry is expected to break when:
(i) multiple branches contribute comparably (mode mixing),
(ii) dispersion becomes significant at shorter wavelengths, or
(iii) geometric/material nonlinearity becomes important at larger amplitudes.
In the present theory these deviations are not treated as flaws; they are structural consequences of a substrate whose timelike direction is selected by the brane action rather than imposed as a metric primitive, and they therefore delimit the domain where a relativistic effective description is expected to apply.
